\documentclass{article}

\usepackage{fancyhdr}
\usepackage{extramarks}

%amsmath
\usepackage{amsmath}
\usepackage{amsthm}
\usepackage{amsfonts}
\usepackage{mathtools}
\usepackage{amssymb}

%Enumerating items and adding columns
\usepackage{enumitem}
\usepackage{multicol}

%TikZ picture
\usepackage{tikz}
\usepackage{scalerel}
\usepackage{pict2e}
\usepackage{tkz-euclide}
\usetikzlibrary{calc}
\usetikzlibrary{patterns,arrows.meta}
\usetikzlibrary{shadows}
\usetikzlibrary{external}
\usetikzlibrary{automata,positioning}

%pgfplots
\usepackage{pgfplots}
\pgfplotsset{compat=newest}
\usepgfplotslibrary{statistics}
\usepgfplotslibrary{fillbetween}

%colours
\usepackage{xcolor}

%
% Basic Document Settings
%

\topmargin=-0.45in
\evensidemargin=0in
\oddsidemargin=0in
\textwidth=6.5in
\textheight=9.0in
\headsep=0.25in

\linespread{1.1}

\pagestyle{fancy}
\lhead{\hmwkAuthorName}
\chead{\hmwkChapter . \hmwkChapterTitle \hmwkClassTime: \hmwkTitle}
\rhead{\firstxmark}
\lfoot{\lastxmark}
\cfoot{\thepage}

\renewcommand\headrulewidth{0.4pt}
\renewcommand\footrulewidth{0.4pt}

\setlength\parindent{0pt}

%
% Create Problem Sections
%

\newcommand{\enterProblemHeader}[1]{
    \nobreak\extramarks{}{Problem \arabic{#1} continued on next page\ldots}\nobreak{}
    \nobreak\extramarks{Problem \arabic{#1} (continued)}{Problem \arabic{#1} continued on next page\ldots}\nobreak{}
}

\newcommand{\exitProblemHeader}[1]{
    \nobreak\extramarks{Problem \arabic{#1} (continued)}{Problem \arabic{#1} continued on next page\ldots}\nobreak{}
    \stepcounter{#1}
    \nobreak\extramarks{Problem \arabic{#1}}{}\nobreak{}
}

\setcounter{secnumdepth}{0}
\newcounter{partCounter}
\newcounter{homeworkProblemCounter}
\setcounter{homeworkProblemCounter}{1}
\nobreak\extramarks{Problem \arabic{homeworkProblemCounter}}{}\nobreak{}

%
% Homework Problem Environment
%
% This environment takes an optional argument. When given, it will adjust the
% problem counter. This is useful for when the problems given for your
% assignment aren't sequential. See the last 3 problems of this template for an
% example.
%
\newenvironment{homeworkProblem}[1][-1]{
    \ifnum#1>0
        \setcounter{homeworkProblemCounter}{#1}
    \fi
    \section{Problem \arabic{homeworkProblemCounter}}
    \setcounter{partCounter}{1}
    \enterProblemHeader{homeworkProblemCounter}
}{
    \exitProblemHeader{homeworkProblemCounter}
}

%
% Homework Details
%   - Title
%   - Due date
%   - Class
%   - Section/Time
%   - Instructor
%   - Author
%

\newcommand{\hmwkTitle}{Exercises}
\newcommand{\hmwkChapter}{CHAPTER 3}
\newcommand{\hmwkDate}{April 3, 2026}
\newcommand{\hmwkClassTime}{}
\newcommand{\hmwkChapterTitle}{Sets}
\newcommand{\hmwkAuthorName}{\textbf{Christian Jelo R. Artoza}}

%
% Title Page
%

\title{
    \vspace{2in}
    \textmd{\textbf{\hmwkChapter:\ \hmwkTitle}}\\
    \normalsize\vspace{0.1in}\small{\hmwkDate}\\
    \vspace{0.1in}\large{\textit{\hmwkChapterTitle\ \hmwkClassTime}}
    \vspace{3in}
}

\author{\hmwkAuthorName}
\date{}

\renewcommand{\part}[1]{\textbf{\large Part \Alph{partCounter}}\stepcounter{partCounter}\\}

%
% Various Helper Commands
%

% qedsymbol that is always flushed to the right
\newcommand{\qedright}{
    \begin{flushright}
        \qedsymbol
    \end{flushright}
}

% For propositions
\newtheorem{prop}{Proposition}

% Useful for algorithms
\newcommand{\alg}[1]{\textsc{\bfseries \footnotesize #1}}

% For derivatives
\newcommand{\deriv}[1]{\frac{\mathrm{d}}{\mathrm{d}x} (#1)}

% For partial derivatives
\newcommand{\pderiv}[2]{\frac{\partial}{\partial #1} (#2)}

% Integral dx
\newcommand{\dx}{\mathrm{d}x}

% Alias for the Solution section header
\newcommand{\solution}{\textbf{\large Solution}}

% Probability commands: Expectation, Variance, Covariance, Bias
\newcommand{\E}{\mathrm{E}}
\newcommand{\Var}{\mathrm{Var}}
\newcommand{\Cov}{\mathrm{Cov}}
\newcommand{\Bias}{\mathrm{Bias}}

\begin{document}

\maketitle

\pagebreak

%Plot Settings
\pgfplotsset{
    standard/.style={
    axis line style = thick,
    axis equal,
    trig format=rad,
    enlargelimits,
    axis x line=middle,
    axis y line=middle,
    axis line style={latex-latex},
    enlarge x limits=0.15,
    enlarge y limits=0.15,
    every axis x label/.style={at={(current axis.right of origin)},anchor=north west},
    every axis y label/.style={at={(current axis.above origin)},anchor=south east}
    }
}


\begin{homeworkProblem}
    \textit{(Exercise 3.7.)} Sketch the following sets as points/arcs/regions in the xy-plane.
    \begin{multicols}{2}
        \begin{enumerate}[label=(\alph*)]
            \item $\{(x,y):x\in [1,3]$ and $y\in[2,4]\}$
            \item $\{(x,y):x\in \mathbb{Z}$ and $y\in \mathbb{R}\}$
            \item $\{(x,y):x^2+y^2=4\}$
            \columnbreak
            \item $\{(x,y):x^2+y^2\leq 4\}$
            \item $\{1,2,3\}\times \{-1,1\}$
            \item $[-1,2]\times [1,\pi]$
        \end{enumerate}
    \end{multicols}
    \textbf{Solution.}
    \begin{enumerate}[label=(\alph*)]
        \item It is the closed rectangle with vertices at $(1,2),(1,4),(3,4),$ and $(3,2)$:
        \begin{center}
            \begin{tikzpicture}
                \begin{axis}[standard,
                    xtick={1,3},
                    ytick={2,4},
                    xlabel={$x$},
                    ylabel={$y$},
                    samples=1000,
                    xmin=-1, xmax=4,
                    ymin=-1, ymax=5
                    ]
                    
                    \node[anchor=center, label=south west:$O$] at (axis cs:0,0){};

                    \draw[blue] (axis cs: 1,2) rectangle (axis cs: 3,4);
                    \fill[blue, fill opacity=0.2] (axis cs: 1,2) rectangle (axis cs: 3,4);
                \end{axis}
            \end{tikzpicture} 
        \end{center}

        \item All vertical lines passing through all integers:
        \begin{center}
            \begin{tikzpicture}
                \begin{axis}[standard,
                    xtick={-2,-1,1,2,3,4},
                    ytick=\empty,
                    xticklabels={$-2$,$-1$,$1$,$2$,$3$,$4$},
                    xticklabel style={anchor=east},
                    xticklabel shift=5pt,
                    xlabel={$x$},
                    ylabel={$y$},
                    samples=1000,
                    xmin=-2, xmax=5,
                    ymin=-1, ymax=5
                    ]
                    
                    \node[anchor=center, label=south west:$O$] at (axis cs:0,0){};

                    \draw[<->,color=blue,] (axis cs:-2,-1.9) -- (axis cs:-2,5.9);
                    \draw[<->,color=blue,] (axis cs:-1,-1.9) -- (axis cs:-1,5.9);
                    \draw[<->,color=blue,] (axis cs:0,-1.9) -- (axis cs:0,5.9);
                    \draw[<->,color=blue,] (axis cs:1,-1.9) -- (axis cs:1,5.9);
                    \draw[<->,color=blue,] (axis cs:2,-1.9) -- (axis cs:2,5.9);
                    \draw[<->,color=blue,] (axis cs:3,-1.9) -- (axis cs:3,5.9);
                    \draw[<->,color=blue,] (axis cs:4,-1.9) -- (axis cs:4,5.9);
                \end{axis}
            \end{tikzpicture} 
        \end{center}

        \item The circle of radius $2$ centered at the origin:
        \begin{center}
            \begin{tikzpicture}
                \begin{axis}[standard,
                    xtick={-2,2},
                    ytick={-2,2},
                    xlabel={$x$},
                    ylabel={$y$},
                    samples=1000,
                    xmin=-2, xmax=2,
                    ymin=-2, ymax=2
                    ]
                    
                    \node[anchor=center, label=south west:$O$] at (axis cs:0,0){};

                    \draw[blue] (0,0) circle (2);
                \end{axis}
            \end{tikzpicture} 
        \end{center}

        \item The closed ball (a.k.a. shaded circle) of radius $2$ centered at the origin:
        \begin{center}
            \begin{tikzpicture}
                \begin{axis}[standard,
                    xtick={-2,2},
                    ytick={-2,2},
                    xlabel={$x$},
                    ylabel={$y$},
                    samples=1000,
                    xmin=-2, xmax=2,
                    ymin=-2, ymax=2
                    ]
                    
                    \node[anchor=center, label=south west:$O$] at (axis cs:0,0){};

                    \draw[blue] (0,0) circle (2);
                    \fill[blue, fill opacity = 0.2] (0,0) circle (2);
                \end{axis}
            \end{tikzpicture} 
        \end{center}

        \item No name; just some points lying around...
            \begin{center}
            \begin{tikzpicture}
                \begin{axis}[standard,
                    xtick={1,2,3},
                    ytick={-1,1},
                    xlabel={$x$},
                    ylabel={$y$},
                    samples=1000,
                    xmin=-1, xmax=3,
                    ymin=-1, ymax=1
                    ]
                    
                    \node[anchor=center, label=south west:$O$] at (axis cs:0,0){};

                    \draw[blue] (1,1) circle (0.05);
                    \fill[blue] (1,1) circle (0.05);

                    \draw[blue] (1,-1) circle (0.05);
                    \fill[blue] (1,-1) circle (0.05);

                    \draw[blue] (2,1) circle (0.05);
                    \fill[blue] (2,1) circle (0.05);

                    \draw[blue] (2,-1) circle (0.05);
                    \fill[blue] (2,-1) circle (0.05);

                    \draw[blue] (3,1) circle (0.05);
                    \fill[blue] (3,1) circle (0.05);

                    \draw[blue] (3,-1) circle (0.05);
                    \fill[blue] (3,-1) circle (0.05);
                \end{axis}
            \end{tikzpicture} 
        \end{center}

        \item The closed rectangle with vertices $(-1,1),(1,\pi),(2,\pi),$ and $(2,1)$:
                \begin{center}
            \begin{tikzpicture}
                \begin{axis}[standard,
                    xtick={-1,2},
                    ytick={1,pi},
                    xticklabels={$-1$,$2$},
                    yticklabels={$1$,$\pi$},
                    xlabel={$x$},
                    ylabel={$y$},
                    samples=1000,
                    xmin=-2, xmax=3,
                    ymin=-1, ymax=4
                    ]
                    
                    \node[anchor=center, label=south west:$O$] at (axis cs:0,0){};

                    \draw[blue, fill opacity=0.2] (-1,1) rectangle (2,pi);
                    \fill[blue, fill opacity=0.2] (-1,1) rectangle (2,pi);
                \end{axis}
            \end{tikzpicture} 
        \end{center}
    \end{enumerate}
    \qedright
\end{homeworkProblem}

\begin{homeworkProblem}
    \textit{(Exercise 3.8.)} Determine whether each of the following is true or false.
    \begin{multicols}{3}
        \begin{enumerate}[label=(\alph*)]
            \item $1\in \{1,\{1\}\}$
            \item $1\subseteq \{1,\{1\}\}$
            \item $1\in \mathcal{P}(\{1,\{1\}\})$
            \item $\{1\}\in \{1,\{1\}\}$
            \item $\{1\}\subseteq \{1,\{1\}\}$
            \columnbreak
            \item $\{1\}\in \mathcal{P}(\{1,\{1\}\})$
            \item $\{\{1\}\}\in\{1,\{1\}\}$
            \item $\{\{1\}\}\subseteq\{1,\{1\}\}$
            \item $\{\{1\}\}\in \mathcal{P}(\{1,\{1\}\})$
            \item $\emptyset \in \mathbb{N}$
            \columnbreak
            \item $\emptyset \subseteq \mathbb{N}$
            \item $\emptyset \in \mathcal{P}(\mathbb{N})$
            \item $\mathbb{Q}\times\mathbb{Q}\subseteq\mathbb{R}\times\mathbb{R}$
            \item $\mathbb{R}^2 \subseteq \mathbb{R}^3$
            \item $\emptyset\subseteq\{1,2,3\}\times\{a,b\}$
        \end{enumerate}
    \end{multicols}
    \textbf{Solution.}
    \begin{multicols}{3}
        \begin{enumerate}[label=(\alph*)]
            \item True
            \item False
            \item False
            \item True
            \item True
            \columnbreak
            \item True
            \item False
            \item True
            \item True
            \item False
            \columnbreak
            \item True
            \item True
            \item True
            \item True
            \item True
        \end{enumerate}
    \end{multicols}
\end{homeworkProblem}

\begin{homeworkProblem}
    \textit{(Exercise 3.9.)} Write down all subsets of each of the following.
    \begin{multicols}{2}
        \begin{enumerate}[label=(\alph*)]
            \item $\{1,2,3\}$
            \item $\{\mathbb{N},\mathbb{Q},\mathbb{R}\}$
            \item $\{\mathbb{N},\{\mathbb{Q},\mathbb{R}\}\}$
            \item $\emptyset$
        \end{enumerate}
    \end{multicols}
    \textbf{Solution.}
    \begin{enumerate}[label=(\alph*)]
        \item $\emptyset,\{1\},\{2\},\{3\},\{1,2\},\{1,3\},\{2,3\},\{1,2,3\}$
        \item $\emptyset,\{\mathbb{N}\},\{\mathbb{Q}\},\{\mathbb{R}\},\{\mathbb{N},\mathbb{Q}\},
        \{\mathbb{N},\mathbb{R}\},\{\mathbb{Q},\mathbb{R}\},\{\mathbb{N},\{\mathbb{Q},\mathbb{R}\}\}$
        \item $\emptyset,\{\mathbb{N}\},\{\{\mathbb{Q},\mathbb{R}\}\},\{\mathbb{N},\{\mathbb{Q},\mathbb{R}\}\}$
        \item $\emptyset$
    \end{enumerate}
\end{homeworkProblem}

\begin{homeworkProblem}
    \textit{(Exercise 3.10.)} The set $\{5a+3b:a,b\in\mathbb{Z}\}$ is equal to a familiar set. By
    examining which elements are possible, determine the familiar set.

    \begin{prop}
        $\{5a+3b:a,b\in\mathbb{Z}\} = \mathbb{Z}$
    \end{prop}

    \begin{proof}
        Said differently, we want to prove that $\{5a+3b:a,b\in\mathbb{Z}\} \subseteq \mathbb{Z}$
        and $\mathbb{Z} \subseteq \{5a+3b:a,b\in\mathbb{Z}\}$. It is clear that the first one is
        true. Indeed, since $a$ and $b$ are integers, it follows that $5a+3b$ is also an integer.
        Now, suppose $n\in \mathbb{Z}$. Then, take $a=-n$ and $b=2n$. And, we get
        \begin{align*}
            5a+3b &= 5(-n)+3(2n) \\
            &= -5n+6n \\
            &= n.
        \end{align*}
        We have shown that, given any $n\in\mathbb{Z}$ there exists $a$ and $b$ such that $n=5a+3b$.
        This means that $\mathbb{Z} \subseteq \{5a+3b:a,b\in\mathbb{Z}\}$, and therefore
        $\{5a+3b:a,b\in\mathbb{Z}\} = \mathbb{Z}$.
    \end{proof}
\end{homeworkProblem}

\begin{homeworkProblem}
    \textit{(Exercise 3.11.)} Suppose $A,B$ and $C$ are sets. Is there a difference between
    $(A\times B)\times C$ and $A\times (B\times C)$? Explain your answer.
    \vspace{1.5mm}

    \textbf{Solution.} They are the same set since they are defined the same.
    \vspace{1.5mm}

    $(A\times B)\times C$ is obtained by first taking the cartesian product of $A$ and $B$,
    which means taking all possible couples $(a,b)$ where $a\in A$ and $b\in B$ to form the
    set $A\times B$, and then taking the cartesian product of the result with $C$. That is,
    combining all $(a,b)$ and all $c$ to form triples $(a,b,c)$.
    \vspace{1.5mm}

    $A\times (B\times C)$ can be described the same way, except we first take the cartesian
    product of $B\times C$ and then take the cartesian product of $A$ and the result.
    \vspace{1.5mm}
    
    In any case, we end up with the same collection of triples $(a,b,c)$ where $a\in A$,
    $b\in B$, and $c\in C$.
    \qedright
\end{homeworkProblem}
\end{document}